\documentclass[12pt,a4paper]{article}

% ── arXiv-compatible packages ──
\usepackage[utf8]{inputenc}
\usepackage[T1]{fontenc}
\usepackage{amsmath,amssymb,amsthm}
\usepackage{hyperref}
\usepackage{geometry}
\usepackage{enumitem}
\usepackage{booktabs}
\geometry{margin=1in}

% ── Theorems ──
\newtheorem{theorem}{Theorem}[section]
\newtheorem{lemma}[theorem]{Lemma}
\newtheorem{corollary}[theorem]{Corollary}
\newtheorem{proposition}[theorem]{Proposition}
\theoremstyle{definition}
\newtheorem{definition}[theorem]{Definition}
\newtheorem{example}[theorem]{Example}
\theoremstyle{remark}
\newtheorem{remark}[theorem]{Remark}

% ── Title ──
\title{\textbf{The Silent Radix: Positional Notation as Ultrametric Tree\\and the Calculus of Indications as Remedy}}

\author{QNFO Research\thanks{Correspondence: research@qnfo.org}}
\date{June 30, 2026}

\begin{document}
\maketitle

\begin{abstract}
Positional notation is not merely a convention for writing numbers---it is an ultrametric tree whose place-value columns encode nested cycles of counting. The radix (base), when made explicit, captures a chosen grouping rhythm that structures the numeral as a hierarchy of nested distinctions. The systematic historical errors catalogued here---the C octal bug, the Mars Climate Orbiter, the Likert-scale mean, the Babylonian sexagesimal ambiguity---all arise from a single structural fault: the substitution of a silent, monocultural decimal default for the domain-specific cyclic grouping, followed by the dissolution of the resulting tree into an Archimedean line. This paper traces the genealogy of this error from Babylonian positional notation through Descartes' analytic geometry to contemporary digital computation. It demonstrates that the $p$-adic numbers are not an exotic alternative but a recovery of the ultrametric already native to positional notation. Re-founding quantitative representation on Spencer-Brown's calculus of indications---where drawing a distinction is marking a cycle---restores the base as an explicit phenomenal frame, the ultrametric as the intrinsic geometry of nested distinctions, zero as the unmarked root, and the self-referential ``10'' as a stable re-entrant form: the cycle measuring its own completion, which is the minimal observer and the formal origin of self-aware measurement.
\end{abstract}

\section{Introduction: The ``10'' Misnomer as Primal Scene}

In any integer base $b \geq 2$, the digit string ``10'' denotes the number $b$. This holds universally: binary ``10'' = two, decimal ``10'' = ten, sexagesimal ``10'' = sixty. Consequently, every base system, when named using its own numeral system, is called ``base-10.'' The phrase ``base-10'' is a misnomer only from an external viewpoint; internally it is perfectly consistent.

This observation---often presented as the joke ``There are 10 types of people in the world: those who understand binary and those who don't'' \cite{hofstadter1979}---is not a mere pun. It exposes a fundamental representational limit: a numeral string cannot disambiguate its own base. The radix is a silent parameter, supplied by the interpretative context---a meta-language, a human convention, a hardware specification.

This paper argues that this ``silent radix'' is not a minor design flaw but the signature of a structural law with consequences spanning computing, science, cognition, and the foundations of mathematics. We trace these consequences through a catalog of 50 documented failures (the Consequence Atlas), demonstrate that the silent radix has a twin---the silent metric of the assumed Euclidean number line---and propose a remedy rooted in Spencer-Brown's \emph{Laws of Form} \cite{spencerbrown1969}.

\section{The Silent Radix in Practice: A Catalog of Consequences}

The Consequence Atlas (accompanying this paper) documents 50 verified failures where a hidden interpretive frame---radix, metric, unit, or scale type---caused error. The taxonomy reveals a consistent pattern.

\subsection{Computing and Software Engineering}

The C programming language's leading-zero-octal convention \cite{ritchie1993} is the paradigmatic silent radix error. The literal \texttt{010} is parsed as octal (decimal 8) rather than decimal 10 because the leading zero---originally inherited from BCPL as a pragmatic notation for machine addresses---serves as an implicit radix marker. This has produced hundreds of documented Common Vulnerabilities and Exposures (CVE) entries, including authentication bypasses where permission masks were silently reinterpreted.

JSON (RFC 8259) mandates decimal-only number representation. Any number computed in hexadecimal, octal, or binary loses its radix on serialization---information is destroyed without warning. The YAML 1.1 ``Norway problem,'' where the country code \texttt{NO} is parsed as boolean \texttt{false} (type coercion), is structurally identical: a silent interpretive rule applied to an unannotated token.

Other documented instances include: JavaScript's \texttt{parseInt("08")} returning \texttt{0} when the radix parameter is omitted; Excel stripping leading zeros from identifiers treated as numbers; IPv4 address parsing where \texttt{010.0.0.1} becomes \texttt{8.0.0.1}; and embedded systems where binary sensor data is displayed as decimal due to missing BCD/radix metadata.

\subsection{Science and Statistics}

The silent metric---the assumption that the number line is straight and equally spaced---produces an analogous class of errors. The most pervasive is the treatment of Likert-scale survey data (ordinal) as interval-level data suitable for means, $t$-tests, and ANOVA. Jamieson \cite{jamieson2004} documented decades of published research using parametric tests on ordinal data without justification.

This error is structurally identical to the silent radix error: a parameter (the metric / scale type) that determines the validity of operations (means vs.\ medians, Pearson $r$ vs.\ Spearman $\rho$) is left implicit and then defaulted to a culturally convenient but inappropriate value.

Clinical pain scales rated 0--10 on a Numerical Rating Scale (NRS) are analyzed parametrically despite pain perception following Steven's power law with exponent $\sim\!3.5$---profoundly nonlinear. The integers 0--10 import the silent metric of equal spacing onto a domain where equal spacing is empirically absent.

\subsection{Physical Units: The Most Catastrophic Frame}

The silent unit is the most consistently catastrophic frame type. Every documented unit error in the Atlas is Critical or Catastrophic:

\begin{itemize}[nosep]
    \item \textbf{Mars Climate Orbiter} (1999, \$327.6M loss): One engineering team used pound-force seconds (imperial) for thruster impulse data while the navigation software expected newton-seconds (metric). The number was transmitted without unit metadata. NASA's mishap investigation board identified ``the failure to use metric units in the coding of a ground software file'' as root cause.
    \item \textbf{Gimli Glider} (1983, 69 lives at risk): Fuel loaded in pounds instead of kilograms during Canada's metric transition. The aircraft ran out of fuel mid-flight.
    \item \textbf{Hubble Space Telescope} (1990, \$1.5B + \$86M corrective optics): Mirror ground to wrong shape because a null corrector spacing error of 1.3mm was not caught. The measurement numbers were ``correct'' within their own frame---but the frame was wrong.
\end{itemize}

\subsection{The Universal Failure Pattern}

Every entry in the Atlas follows the same structural sequence:

\begin{enumerate}[nosep]
    \item A quantitative representation is generated in one context with a specific interpretive frame (radix, metric, unit, scale type).
    \item The frame is not carried with the representation.
    \item The representation enters a new context where the default frame differs.
    \item The mismatch produces an error.
\end{enumerate}

This pattern is invariant across domains. A C compiler parsing \texttt{010}, a NASA navigation team receiving unlabeled thruster data, a researcher computing means on Likert data, and a Babylonian scribe reading ambiguous cuneiform---all fell into the same structural trap.

\section{The Genealogy of the Silent Frame}

\subsection{Pre-Positional Systems: Transparency Without Scale}

Tally sticks and token systems (Upper Paleolithic through Sumerian, $\sim$8000 BCE) encoded quantity by one-to-one correspondence. There was no base, no positional value---and therefore no silent frame. The limitation was obvious: the system did not scale. Additive numeral systems (Egyptian hieroglyphs, Roman numerals) used explicit symbols for powers of the base: \texttt{X} always meant ten, regardless of position. These systems avoided the silent radix at the cost of algorithmic inefficiency---a trade-off documented by the persistence of Roman numerals in European accounting until the 14th century precisely because their explicit values made forgery harder.

\subsection{The Positional Revolution: Power Through Silence}

Babylonian sexagesimal notation ($\sim$2000 BCE) introduced place value but lacked a consistent zero placeholder. The digit \texttt{1} could mean 1, 60, 3600, or 1/60 depending on context. This was the original silent radix---positional notation's power was purchased at the cost of interpretive ambiguity. The Indian invention of zero as a placeholder ($\sim$5th--7th century CE) resolved the place-value ambiguity but introduced a new silence: the radix became decimal by default, invisible in the notation itself.

\subsection{The Flattening: From Tree to Line}

Simon Stevin's \emph{De Thiende} (1585) extended decimal positional notation to fractions, and Ren\'e Descartes' \emph{La G\'eom\'etrie} (1637) identified numbers with points on a geometric line. This unified the discrete tree of positional notation with the continuous, equally spaced Euclidean line. The 19th-century arithmetization of analysis (Cauchy, Weierstrass, Dedekind) canonized $\mathbb{R}$ as \emph{the} complete ordered field---obscuring the existence of alternative completions.

The silent metric was thus sealed: the tree became a line, and the line became the fundamental reality of number, rather than a derived abstraction built on a discrete, ultrametric scaffolding.

\subsection{Digital Recapitulation}

Modern computing inherited this entire layered history. Programming languages embedded silent radices (C's leading-zero octal), silent units (floating-point without unit annotation), and silent scale types (integers treated as interval by default). Automation locked these defaults into infrastructure, making them physically enforced and resistant to correction.

\section{The Mathematical Demonstration: $p$-adic Numbers Recover the Native Ultrametric}

\subsection{Positional Notation Is Already Ultrametric}

Positional notation inherently carries an ultrametric structure. In base $b$, the distance between two integers measured by the highest power of $b$ dividing their difference:
\begin{equation}
|x - y|_b = b^{-v_b(x-y)}
\end{equation}
where $v_b(n)$ is the exponent of the highest power of $b$ dividing $n$. This is exactly the $b$-adic valuation---a non-Archimedean (ultrametric) absolute value. Two numbers that share more trailing digits (more right-aligned positional agreement) are closer in this metric.

The place-value columns create a hierarchical branching structure: a tree, not a line. This tree is the native geometry of positional notation---the ultrametric is not an interpretation added later but a structural property of the digit-string representation.

\subsection{$p$-adic Numbers as Generalization}

For prime $p$, the $p$-adic numbers $\mathbb{Q}_p$ are the completion of $\mathbb{Q}$ under the $p$-adic absolute value. Topologically, $\mathbb{Z}_p$ is a Cantor set---a tree where every ball is partitioned into $p$ disjoint sub-balls. The strong triangle inequality holds:
\begin{equation}
|x - z|_p \leq \max(|x - y|_p, |y - z|_p)
\end{equation}

This produces the property that all triangles are isosceles with the base at most the legs, and every point in a disc is its center. There is no total order compatible with the topology.

\subsection{Ostrowski's Theorem and the Contingency of the Line}

Ostrowski's Theorem (1916) \cite{ostrowski1916} states that the only non-trivial absolute values on $\mathbb{Q}$ are the Euclidean one and the $p$-adic ones. This proves that the Euclidean metric of the number line is one choice among infinitely many, each equally forced by arithmetic. The decimal number line---the Archimedean completion---is not the unique truth about quantity; it is a culturally reinforced default.

The $p$-adic numbers demonstrate that it is possible to have a number system where the base and metric are intrinsic and visible. The base $p$ is constitutive of the number system; the metric is algebraically determined by divisibility, not imported as a geometric convention.

The ``errors'' of the silent radix and silent metric are instances of projecting the wrong completion onto a representation whose native structure is ultrametric. The $p$-adic view recovers what was lost.

\section{The Epistemological Core: Ontology, Epistemology, and the Gap}

\subsection{The Use/Mention Boundary}

The ``10'' problem is a miniature model of the use--mention distinction \cite{quine1940} and of Tarski's undefinability of truth \cite{tarski1933}. Inside a given base system, ``10'' is ontologically the base---a complete object with properties. But across systems, ``10'' shifts referent---the same string denotes different numbers. The base cannot be determined from the string alone; the system cannot ground itself.

This is the same structural limit that produces G\"odel's incompleteness \cite{godel1931}: any sufficiently expressive formal system requires an external meta-language to resolve self-referential statements. The silent radix is a G\"odel sentence in miniature---a concrete, operational undecidability at the level of notation, not just arithmetic.

\subsection{The Collapse and Its Consequences}

In any functioning symbolic system, ontology (what a thing is) and epistemology (how it is known) are operationally indistinguishable. However, the act of defining or changing the system requires an external vantage---a meta-language. The gap between ontology and epistemology is not a defect to be eliminated but the condition of meaning.

The silent-frame error is precisely the mistake of collapsing this gap: of taking the internal ontology (``10 means ten, the line is straight'') as absolute, forgetting the epistemic act that constructed it. Na\"ive Platonism treats numbers as purely ontological; na\"ive idealism treats them as purely epistemological. The responsible stance is to hold both together, with the meta-awareness that the frame is always a choice.

\section{The Remedy: Laws of Form as a Foundation of Explicit Distinction}

\subsection{Set Theory's Hidden Assumptions}

Standard Zermelo-Fraenkel set theory (ZFC) builds everything from the empty set and membership ($\in$). Membership is an unexamined primitive; the act of distinguishing the empty set is erased. Numbers are constructed as extensional sets (von Neumann ordinals), in which neither the base nor the metric plays any foundational role. The observer---the one who draws the distinction---is systematically excluded.

This foundational amnesia is the root cause of the silent-frame problem. If the ground of mathematics erases the act of distinction, then all representations built on that ground will systematically lose their interpretive frames.

\subsection{The Calculus of Indications}

Spencer-Brown's \emph{Laws of Form} \cite{spencerbrown1969} offers a radical alternative. Its primitive is not membership or the empty set, but the act of drawing a distinction. A distinction cleaves a space into a marked state and an unmarked state. The two primitive laws---calling (a distinction repeated is the same distinction) and crossing (crossing a distinction twice returns to the unmarked state)---generate a complete calculus.

In this framework:

\begin{itemize}[nosep]
    \item \textbf{Zero is the unmarked state}, the root of all counting. All numbers grow from this root by adding distinctions.
    \item \textbf{A numeral is a nested form} whose depth structure visibly encodes the base. The base is the arity of distinction at each level---a visible parameter of the form itself.
    \item \textbf{The ultrametric is the geometry of the form}: two numbers are close to the depth of their shared nesting---the number of matching trailing distinctions. This is not an added metric but intrinsic to the construction.
    \item \textbf{Re-entry} allows a form to enter its own space, producing self-reference without paradox. The re-entrant form $f = \neg f$ oscillates, generating time and memory within the calculus.
\end{itemize}

\subsection{``10'' as Stable Re-entrant Form}

The self-referential ``10''---the base naming itself---is precisely a re-entrant form. In a tree of cycles, ``10'' means: one cycle at the next level, zero cycles at the current level. This is the completion of the finest level (the inner space) crossing the boundary and appearing as a mark at the next level. The string ``10'' is the boundary crossing made visible.

In LoF, this is not paradoxical. It is a stable oscillation: the cycle that marks its own completion. This is the minimal observer---the act of registering that a cycle has completed and thereby creating the next level. The silent radix error is the forgetting of this observer. Making the base explicit restores the observer to the representation.

\subsection{Recovery of the Euclidean Line}

The Archimedean line is not rejected. It is recovered as a secondary projection---an explicit completion of the discrete tree under a declared metric. In this layered model, the tree is primary; the line is derived and flagged. The error was never to use the line; it was to forget the tree from which the line was grown.

\section{The Thesis Stated}

\textbf{Positional notation is a rooted tree of nested time-cycles.} Its radix---when explicitly chosen to match phenomenal periodicities---inscribes the ultrametric of shared temporal nesting into the numeral. The historical errors of silent radices, silent metrics, and assumed linearity are consequences of replacing this grounded, cyclic tree with a monocultural, unmarked decimal default and then flattening it into an Archimedean line, erasing both the time-structure of the measured world and the observer whose act of distinction generates the tree. Re-founding on the calculus of indications restores the base as an explicit phenomenal frame, the ultrametric as the native geometry of nested time, zero as the unmarked root, and the self-referential ``10'' as the stable re-entrant form of the cycle counting its own completion---the minimal observer, the seed of all self-aware measurement, and the formal origin of time in the act of distinction.

\section{Nine Principles for a Cyclic-Frame Quantitative Practice}

\begin{enumerate}[nosep]
    \item \textbf{Cyclic Grounding.} Every numeral shall declare the cycles it counts. The base corresponds to observed or chosen periodicities.
    \item \textbf{Explicit Frame.} No number shall be transmitted, stored, or operated upon without its radix, metric, unit, and scale type.
    \item \textbf{Native Ultrametric.} The default distance is shared nesting depth. Linear distance is an applied, flagged transformation.
    \item \textbf{Zero as Root.} The unmarked state is the origin. There is no ``negative distance'' from the root.
    \item \textbf{Re-entry as Self-Measurement.} ``10'' is the unit cycle observing its own completion---the formal atom of reflexivity.
    \item \textbf{Arithmetic Invariance.} The laws hold across all frames, but phenomenal meaning is frame-dependent.
    \item \textbf{Layered Completion.} The continuous is a derived abstraction on the discrete tree. The tree is primary; the line is secondary and flagged.
    \item \textbf{Second-Order Numeracy.} Any system handling numbers must represent its own representational choices, or it is first-order and vulnerable to silent-frame error.
    \item \textbf{Temporal Primacy.} Time, as the counting of cycles, is the fundamental quantity. Space, as the continuous line, is a derived flattening.
\end{enumerate}

\section{Conclusion: From Misnomer to Foundation}

The conversation that began with a joke about ``10'' ends with a new foundation for quantitative representation. The silent radix is not a quirk to be tolerated but the signature of a structural law: representation requires interpretation, and interpretation requires an explicit frame. When the frame is hidden---in a C literal, a JSON payload, a Likert scale, or a spacecraft trajectory file---error becomes inevitable.

The calculus of indications offers a foundation in which frames cannot be hidden, because every form is an act of distinction and every distinction carries its boundary visibly. In this foundation, the ``10'' misnomer is transformed from a bug into the engine of self-aware measurement---the tree observing its own branching, the cycle counting its own completion.

% ── arXiv metadata ──
% arxiv: submit to math.HO and cs.LO
\section*{Acknowledgments}
This work is part of the Silent Radix Research Program. The accompanying Consequence Atlas documents 50 verified silent-frame errors across computing, science, statistics, and engineering. The formal mathematical appendix contains eight theorems with proof sketches. All artifacts are available at DOI \href{https://doi.org/10.5281/zenodo.21052039}{10.5281/zenodo.21052039}.

\bibliographystyle{plain}
\begin{thebibliography}{99}

\bibitem{dehaene1997}
S.~Dehaene.
\newblock \emph{The Number Sense}.
\newblock Oxford University Press, 1997.

\bibitem{descartes1637}
R.~Descartes.
\newblock \emph{La G\'eom\'etrie}, 1637.

\bibitem{goodman1968}
N.~Goodman.
\newblock \emph{Languages of Art}.
\newblock Hackett, 1968.

\bibitem{godel1931}
K.~G\"odel.
\newblock \"Uber formal unentscheidbare S\"atze.
\newblock \emph{Monatshefte f\"ur Mathematik und Physik}, 1931.

\bibitem{hensel1897}
K.~Hensel.
\newblock \"Uber eine neue Begr\"undung der Theorie der algebraischen Zahlen.
\newblock \emph{Jahresbericht der DMV}, 1897.

\bibitem{hofstadter1979}
D.~Hofstadter.
\newblock \emph{G\"odel, Escher, Bach}.
\newblock Basic Books, 1979.

\bibitem{jamieson2004}
S.~Jamieson.
\newblock Likert scales: how to (ab)use them.
\newblock \emph{Medical Education}, 2004.

\bibitem{kauffman1987}
L.~Kauffman.
\newblock Arithmetic in the Calculus of Indications.
\newblock \emph{International Journal of General Systems}, 1987.

\bibitem{maturana1980}
H.~Maturana and F.~Varela.
\newblock \emph{Autopoiesis and Cognition}.
\newblock Reidel, 1980.

\bibitem{ostrowski1916}
A.~Ostrowski.
\newblock \"Uber einige L\"osungen der Funktionalgleichung.
\newblock \emph{Acta Mathematica}, 1916.

\bibitem{quine1940}
W.~V.~O.~Quine.
\newblock \emph{Mathematical Logic}.
\newblock Harvard University Press, 1940.

\bibitem{ritchie1993}
D.~Ritchie.
\newblock The Development of the C Language.
\newblock \emph{ACM HOPL-II}, 1993.

\bibitem{schikhof1984}
W.~Schikhof.
\newblock \emph{Ultrametric Calculus}.
\newblock Cambridge University Press, 1984.

\bibitem{siegler2003}
R.~Siegler and J.~Opfer.
\newblock The Development of Numerical Estimation.
\newblock \emph{Psychological Science}, 2003.

\bibitem{smith1996}
B.~C.~Smith.
\newblock \emph{On the Origin of Objects}.
\newblock MIT Press, 1996.

\bibitem{spencerbrown1969}
G.~Spencer-Brown.
\newblock \emph{Laws of Form}.
\newblock Allen \& Unwin, 1969.

\bibitem{stevens1946}
S.~S.~Stevens.
\newblock On the Theory of Scales of Measurement.
\newblock \emph{Science}, 1946.

\bibitem{stevin1585}
S.~Stevin.
\newblock \emph{De Thiende}, 1585.

\bibitem{tarski1933}
A.~Tarski.
\newblock The Concept of Truth in Formalized Languages, 1933.

\bibitem{varela1979}
F.~Varela.
\newblock \emph{Principles of Biological Autonomy}.
\newblock North-Holland, 1979.

\bibitem{vonfoerster1974}
H.~von Foerster.
\newblock \emph{Cybernetics of Cybernetics}, 1974.

\bibitem{wittgenstein1956}
L.~Wittgenstein.
\newblock \emph{Remarks on the Foundations of Mathematics}.
\newblock Blackwell, 1956.

\end{thebibliography}

\end{document}
